\documentclass[a4paper,10pt]{report}\input{../0.Préambule/Préambule_cours_prof}\begin{document}

\definecolor{titlecolor}{rgb}{0.98, 0.4, 0.0}
\newpage
\setcounter{chapter}{2}
\chapter{Quotient de nombres relatifs}

\section{Division de deux nombres relatifs}

\begin{regle}
Pour calculer le quotient d'un nombre relatif par un nombre relatif non nul, on divise leur distance à zéro et on applique la règle des signes suivante :
\begin{enumerate}
\item[•] le quotient de deux nombres relatifs de même signe est positif ;
\item[•] le quotient de deux nombres relatifs de signes contraires est négatif.
\end{enumerate}
\end{regle}

\begin{ex}
Effectue la division suivante :$ K=65 \div (-5)$

\noindent\grandscarreaux{\linewidth}{4cm}
\end{ex}

\begin{ex}
Quelle est l'écriture décimale du quotient $L=\dfrac{(-30)}{(-4)}$ ?

\noindent\grandscarreaux{\linewidth}{4cm}
\end{ex}

\begin{rmq}
\begin{enumerate}
\item[•] La règle des signes pour la division est la même que celle pour la multiplication.
\item[•] Le quotient de $0$ par n'importe quel nombre non nul est égal à $0$.
\end{enumerate}
\end{rmq}

Cela signifie que pour tout nombre relatif non nul $a$, on a $\dfrac{0}{a}=0$

\section{Valeurs approchées}

\begin{dft}
Une valeur approchée est un nombre proche de la valeur exacte de ce nombre. On utilise ces valeurs à la place du véritable nombre pour simplifier des calculs
\end{dft}

\begin{ex}

\noindent\grandscarreaux{\linewidth}{4cm}
\end{ex}

\begin{meth}[troncature ou arrondi ?]

La fraction \href{https://www.numworks.com/fr/simulateur/}{$\dfrac{37}{23}$} ayant un quotient décimal infini, on souhaite ne donner une valeur approchée au millième. 

\begin{center}
\textbf{VALEUR APPROCHEE}
\end{center}

\noindent \begin{minipage}{8cm}
\begin{center}
\textbf{Troncature}
\end{center}
Pour donner la troncature au millième, on coupe le nombre décimal juste après sa troisième décimale :

\vspace{0.9cm}
\begin{center}
\huge{$1,608$} \bcstop~ \textcolor{lightgray}{$695652...$}
\end{center}

\vspace{0.9cm}
La troncature au millième de $\dfrac{37}{13}$ est $1,608$
\end{minipage}
\hspace{3mm}
\vline
\hspace{3mm}
\begin{minipage}{8cm}
\begin{center}
\textbf{Arrondi}
\end{center}
Pour donner l'arrondi au millième, Il faut regarder la quatrième décimale :

\begin{enumerate}
\item[•]Si il s'agit du $0$, $1$, $2$, $3$, $4$ alors l'arrondi est équivalent à la troncature.
\item[•]Si il s'agit du $5$, $6$, $7$, $8$, $9$ alors on augment de $1$ la troisième décimale.
\end{enumerate}
\begin{center}
\huge{$1,608$}\textcolor{red}{$6$}\textcolor{lightgray}{$695652...$}
\end{center}
La quatrième décimale étant un $6$, l'arrondi au millième de $\dfrac{37}{23}$ est donc $1,609$.
\end{minipage}
\end{meth}

\section{Inverse d'un nombre relatif}

\begin{dft}
Deux nombres sont inverses l'un de l'autre signifie que leur produit est égal à $1$.

\noindent Autrement dit, soient $a$ et $b$ deux nombres non nuls, $a$ est l'inverse de $b$ si $a \times b = 1$.
\end{dft}

\begin{ex}
\noindent\grandscarreaux{\linewidth}{3.2cm}
\end{ex}

\begin{rmq}

\vspace{-4.5mm} \hspace{3cm}$0$ n'a pas d'inverse.
\end{rmq}

\begin{prop}
$\bullet$ L'inverse de $a$ est $\dfrac{1}{a}$. \hspace{3cm} $\bullet$ L'inverse de $\dfrac{1}{a}$ est $a$. \hspace{3cm} $\bullet$ L'inverse de $\dfrac{a}{b}$ est $\dfrac{b}{a}$.
\end{prop}
\end{document}